\documentclass[a4paper,12pt]{article}
\usepackage[T2A]{fontenc}
\usepackage[utf8]{inputenc}
\usepackage[russian]{babel}
\usepackage{amsmath,amssymb,mathtools}
\usepackage{geometry}
\geometry{left=23mm,right=22mm,top=21mm,bottom=22mm}
\usepackage{microtype}
\usepackage{icomma}
\usepackage{xcolor}
\definecolor{accent}{HTML}{0B7779}
\definecolor{darkaccent}{HTML}{145254}
\definecolor{textgray}{HTML}{414B4B}
\usepackage{tikz}
\usetikzlibrary{arrows.meta,positioning,shapes.geometric,calc}
\usepackage{tabularx,booktabs,array,longtable}
\usepackage{enumitem}
\usepackage{fancyhdr}
\usepackage{setspace}
\usepackage{hyperref}
\hypersetup{colorlinks=true,linkcolor=darkaccent,urlcolor=darkaccent,pdftitle={Чек-лист: экономическая оптимизация},pdfauthor={Лёвин Артём Александрович}}
\onehalfspacing
\setlength{\parindent}{0pt}
\setlength{\parskip}{5pt plus 1pt minus 1pt}
\setlist[enumerate]{leftmargin=7mm,itemsep=5pt,topsep=3pt}
\setlist[itemize]{leftmargin=6mm,itemsep=3pt,topsep=2pt}
\setcounter{secnumdepth}{2}
\renewcommand{\thesection}{\arabic{section}}
\renewcommand{\thesubsection}{\arabic{section}.\arabic{subsection}}
\pagestyle{fancy}
\fancyhf{}
\fancyhead[L]{\footnotesize\color{textgray} Математика \quad/\quad Экономическая оптимизация}
\fancyhead[R]{\footnotesize\color{textgray} Матвей Горбачев}
\fancyfoot[C]{\footnotesize\color{textgray}\thepage}
\renewcommand{\headrulewidth}{0.3pt}
\setlength{\headheight}{15pt}
\newcommand{\keyidea}[1]{\par\noindent\textcolor{darkaccent}{\textbf{Ключевая идея.}}~#1\par}
\newcommand{\caution}[1]{\par\noindent\textcolor{darkaccent}{\textbf{Проверьте.}}~#1\par}
\newcommand{\worked}[1]{\subsection*{#1}}
\begin{document}
\thispagestyle{empty}
\begin{center}
\vspace*{24mm}
{\large\color{darkaccent} МАТЕМАТИКА\par}
\vspace{13mm}
{\LARGE Чек-лист\par}
\vspace{9mm}
{\huge\color{accent} Экономические задачи\\[2mm] на оптимизацию\par}
\vspace{15mm}
{\large Подготовка к ЕГЭ по профильной математике\par}
\vfill
{\Large Матвей Горбачев\par}
\vspace{10mm}
{\large 15 ноября 2026 г.\par}
\vspace{15mm}
{\normalsize Лёвин Артём Александрович\par}
{\normalsize эксклюзивно для Матвея Горбачева\par}
\vspace{20mm}
\begin{tikzpicture}[x=1cm,y=1cm]
\draw[accent,line width=0.9pt] (0,0) .. controls (1.0,0.55) and (2.5,-0.48) .. (3.5,0)
  .. controls (4.5,0.48) and (6.0,-0.56) .. (7,0);
\end{tikzpicture}
\end{center}
\clearpage

\begingroup\setstretch{1.03}
\section{Главная цель задачи}
В задачах на оптимизацию нужно выбрать \textbf{лучшее допустимое значение}: наименьшие расходы, наименьшую цену либо наибольшее количество произведённой продукции.

\keyidea{Сначала выясните, \emph{какую величину} требуется сделать минимальной или максимальной. Именно для неё составляют целевую функцию.}

\begin{tabularx}{\linewidth}{@{}p{0.30\linewidth}X@{}}
\toprule
\textbf{Формулировка вопроса} & \textbf{Целевая функция}\\
\midrule
«Наименьшая зарплата» & Общая выплата рабочим $S$.\\
«Наименьшие затраты» & Общие расходы $C$.\\
«Наибольшее число изделий» & Общий выпуск $N$.\\
«Наименьшая цена» & Цена единицы продукции $p$.\\
\bottomrule
\end{tabularx}

\section{Как составить математическую модель}
\begin{enumerate}
\item \textbf{Обозначьте независимую переменную.} Например, $t$~--- число рабочих на первом объекте.
\item \textbf{Создайте таблицу.} Для каждого объекта укажите количество рабочих, изделий, часов или денег, в зависимости от текста задачи.
\item \textbf{Используйте связь между переменными.} Если всего $24$ человека, а на первом объекте $t$, то на втором $24-t$. Если всего $20$ изделий, то $t_1+t_2=20$.
\item \textbf{Составьте целевую функцию одной переменной.} Подставьте найденную связь в сумму затрат либо выпуска.
\item \textbf{Зафиксируйте ограничения.} Количество людей или изделий целое и неотрицательное; знаменатель не равен нулю; подкоренное выражение неотрицательно.
\end{enumerate}

\caution{Формула «если изготовлено $u$ изделий, расходы составляют $4u^3$» действует при любом допустимом $u$. Если изготовлено $20-t$ изделий, расходы равны $4(20-t)^3$. Скобки обязательны.}

\section{Исследование функции}
Для целевой функции $F(x)$ найдите производную $F'(x)$, решите $F'(x)=0$ и определите знак производной на допустимых промежутках:
\[
F'(x)<0\ \Rightarrow\ F\text{ убывает},\qquad
F'(x)>0\ \Rightarrow\ F\text{ возрастает}.
\]
Проверьте границы допустимого промежутка и учитывайте, \textbf{должны ли ответы быть целыми}. Непрерывная точка экстремума может оказаться недостижимой в реальной задаче.

\endgroup
\begingroup\setstretch{1.03}\setlength{\parskip}{2pt}
\section{Пример: распределение работников}
\textbf{Условие.} Бригада из $24$ человек распределяется между двумя объектами. Если на первом объекте работают $t$ человек, общая выплата им составляет $4t^2$ условных единиц. Если на втором объекте работают $u$ человек, общая выплата им равна $u^2$ условных единиц. Сколько работников направить на каждый объект, чтобы общая выплата была наименьшей?

\textbf{Таблица модели}
\begin{center}
\begin{tabular}{@{}lcc@{}}
\toprule
Объект & Работников & Выплата, у.~е.\\
\midrule
Первый & $t$ & $4t^2$\\
Второй & $24-t$ & $(24-t)^2$\\
\bottomrule
\end{tabular}
\end{center}

\textbf{Целевая функция и ограничения:}
\[
S(t)=4t^2+(24-t)^2=5t^2-48t+576,
\qquad t\in\{0,1,\ldots,24\}.
\]
\textbf{Непрерывное исследование:}
\[
S'(t)=10t-48=0\ \Rightarrow\ t=4{,}8.
\]
При $t<4{,}8$ производная отрицательна, при $t>4{,}8$ положительна. Значит, непрерывный минимум достигается при $4{,}8$.

\textbf{Учёт целочисленности.} Количество людей целое. Проверяем соседние допустимые варианты:
\begin{align*}
S(4)&=4\cdot4^2+20^2=64+400=464,\\
S(5)&=4\cdot5^2+19^2=100+361=461.
\end{align*}
\textcolor{darkaccent}{\textbf{Ответ:}} на первом объекте $5$ работников, на втором $19$; минимальная выплата~--- $461$~у.~е.

\caution{Значения $4$ и $5$ подставляют в исходную целевую функцию $S$, а для определения возрастания и убывания используют её производную $S'$.}

\section{Пример: два завода и кубическая зависимость}
\textbf{Условие.} Два завода должны изготовить за неделю всего $20$ приборов. На первом заводе изготовление $t_1$ приборов требует $4t_1^3$ человеко-часов, на втором изготовление $t_2$ приборов требует $t_2^3$ человеко-часов. Один человеко-час оплачивается по ставке $1000$ рублей. Как распределить выпуск, чтобы расходы на оплату труда были минимальны?

Из равенства $t_1+t_2=20$ получаем $t_2=20-t_1$. Считаем затраты в тысячах рублей:
\[
C(t_1)=4t_1^3+(20-t_1)^3,\qquad t_1\in\{0,1,\ldots,20\}.
\]
\[
C'(t_1)=12t_1^2-3(20-t_1)^2
=3(2t_1-(20-t_1))(2t_1+(20-t_1)).
\]
\[
C'(t_1)=0\quad\Rightarrow\quad t_1=\frac{20}{3}
\quad\text{или}\quad t_1=-20\ (\text{вне области}).
\]
На допустимом промежутке $C$ убывает до $20/3$ и возрастает после него.\par\pagebreak[4]\noindent Проверяем $6$ и $7$:
\begin{align*}
C(6)&=4\cdot6^3+14^3=3608,\\
C(7)&=4\cdot7^3+13^3=3569.
\end{align*}
\textcolor{darkaccent}{\textbf{Ответ:}} первый завод изготавливает $7$ приборов, второй~--- $13$; расходы~--- $3569$ тысяч рублей.

\section{Пример: максимум выпуска при ограниченном бюджете}
\textbf{Условие.} Два завода выпускают одинаковые изделия. Если первый выпускает $t_1$ единиц, расходы на оплату труда равны $250t_1^2$ рублей; если второй выпускает $t_2$ единиц, расходы равны $200t_2^2$ рублей. Общий бюджет на неделю~--- $900\,000$ рублей. Сколько изделий максимально можно изготовить за неделю?

В оптимальной непрерывной модели весь бюджет можно использовать: увеличение любого объёма выпуска увеличивает выпуск, пока есть свободные средства. Поэтому
\[
250t_1^2+200t_2^2=900\,000.
\]
Отсюда, учитывая $t_2\ge0$,
\[
t_2=\sqrt{4500-1{,}25t_1^2},\qquad 0\le t_1\le60.
\]
\textbf{Теперь максимизируем выпуск:}
\[
N(t_1)=t_1+t_2=t_1+\sqrt{4500-1{,}25t_1^2}.
\]
На внутренней части области определения
\[
N'(t_1)=1-\frac{1{,}25t_1}{\sqrt{4500-1{,}25t_1^2}}.
\]
\[
N'(t_1)=0\quad\Rightarrow\quad
\sqrt{4500-1{,}25t_1^2}=1{,}25t_1.
\]
Обе части неотрицательны, поэтому возведение в квадрат допустимо:
\begin{align*}
4500-1{,}25t_1^2&=1{,}5625t_1^2,\\
4500&=2{,}8125t_1^2,\\
t_1^2&=1600,\qquad t_1=40.
\end{align*}
Тогда $t_2=\sqrt{4500-1{,}25\cdot1600}=50$.

При $0<t_1<40$ производная положительна, при $40<t_1<60$ отрицательна. Следовательно, это глобальный максимум на $[0,60]$; на границах значения меньше. Объёмы $40$ и $50$ целые и допустимые:
\[
250\cdot40^2+200\cdot50^2=900\,000.
\]
\textcolor{darkaccent}{\textbf{Ответ:}} $90$ изделий: $40$ на первом заводе и $50$ на втором.


\section{Пример: наименьшая цена продукции}
\textbf{Условие.} Годовые расходы предприятия (в миллионах рублей) при выпуске $x$ тысяч изделий составляют
\[
q(x)=0{,}5x^2+x+7.
\]
Цена одного изделия равна $p$ тысяч рублей, поэтому годовая выручка равна $px$ миллионам рублей. Производство ежегодно работает по одному и тому же плану. Найдите наименьшее значение $p$, при котором суммарная прибыль за три года составит не менее $75$ млн рублей.

Годовая прибыль равна $px-q(x)$ млн рублей. За три года требуется:
\[
3\bigl(px-q(x)\bigr)\ge75\quad\Rightarrow\quad
px-\bigl(0{,}5x^2+x+7\bigr)\ge25.
\]
При $x>0$ выражаем нижнюю границу цены:
\[
p\ge\frac{25+0{,}5x^2+x+7}{x}
=\underbrace{\frac{x}{2}+1+\frac{32}{x}}_{\displaystyle P(x)}.
\]
Значит, минимальная цена равна минимуму $P(x)$ при $x>0$.
\[
P'(x)=\frac12-\frac{32}{x^2}
=\frac{0{,}5x^2-32}{x^2}.
\]
\[
P'(x)=0\quad\Rightarrow\quad x^2=64\quad\Rightarrow\quad x=8.
\]
При $0<x<8$ производная отрицательна, а при $x>8$ положительна. Следовательно,
\[
P_{\min}=P(8)=4+1+4=9.
\]
\textbf{Проверка по условию:} $q(8)=0{,}5\cdot64+8+7=47$ млн руб.; годовая выручка $px=9\cdot8=72$ млн руб.; годовая прибыль $72-47=25$ млн руб.; за три года получаем $75$ млн руб.

\textcolor{darkaccent}{\textbf{Ответ:}} минимальная цена~--- $9$ тысяч рублей за изделие (при выпуске $8$ тысяч изделий ежегодно).

\endgroup
\clearpage
\thispagestyle{fancy}
\begin{center}
{\Large\color{darkaccent} Алгоритм решения экономической задачи на оптимизацию}
\end{center}
\vspace{7mm}
\begin{center}
\begin{tikzpicture}[
  >=Latex,
  node distance=7.5mm,
  action/.style={draw=accent,rounded corners=2mm,line width=.65pt,align=center,text width=117mm,minimum height=12.5mm,font=\small,inner sep=3pt},
  terminal/.style={draw=darkaccent,ellipse,line width=.75pt,align=center,text width=72mm,minimum height=12mm,font=\small,inner sep=3pt},
  decide/.style={draw=accent,diamond,aspect=3.7,align=center,text width=62mm,inner sep=1.6pt,font=\small,minimum height=17mm},
  branch/.style={draw=accent,rounded corners=2mm,line width=.65pt,text width=49mm,minimum height=18mm,align=center,font=\small,inner sep=3pt},
  arr/.style={->,line width=.8pt,color=darkaccent}
]
\node[terminal] (start) {Прочитайте вопрос: что нужно сделать минимальным или максимальным?};
\node[action,below=of start] (table) {Назовите переменные. Составьте таблицу «объект~--- количество~--- расходы/выпуск».};
\node[action,below=of table] (constraint) {Запишите связь между переменными и все ограничения.};
\node[action,below=of constraint] (func) {Выразите целевую функцию через одну переменную $F(x)$.};
\node[action,below=of func] (deriv) {Исследуйте $F'(x)$, критические точки, знаки и границы области.};
\node[decide,below=12mm of deriv] (integer) {Ответ обязан быть целым?};
\node[branch,below left=15mm and 9mm of integer] (yes) {Сравните допустимые целые варианты и докажите оптимальность.};
\node[branch,below right=15mm and 9mm of integer] (no) {Используйте допустимую точку экстремума и проверьте границы.};
\node[terminal,below=47mm of integer] (end) {Запишите ответ с единицами измерения и проверкой.};
\draw[arr] (start)--(table);
\draw[arr] (table)--(constraint);
\draw[arr] (constraint)--(func);
\draw[arr] (func)--(deriv);
\draw[arr] (deriv)--(integer);
\draw[arr] (integer.west) -| node[pos=.17,above,font=\scriptsize]{да} (yes.north);
\draw[arr] (integer.east) -| node[pos=.17,above,font=\scriptsize]{нет} (no.north);
\draw[arr] (yes.south) |- (end.north west);
\draw[arr] (no.south) |- (end.north east);
\end{tikzpicture}
\end{center}
\clearpage

\section*{Тренировочный блок}
\addcontentsline{toc}{section}{Тренировочный блок}
\textcolor{darkaccent}{\textbf{Путь А. 10 задач от базовых к более сложным.}} Для каждой задачи составьте модель, запишите допустимые значения переменных, исследуйте целевую функцию и дайте ответ с единицами измерения. Для целочисленных величин обязательно проверьте допустимые целые значения.

\subsection*{Базовый уровень}
\begingroup\small\setstretch{1.02}
\setlength{\parskip}{2pt}
\setlist[enumerate]{itemsep=1pt,topsep=0pt}
\begin{enumerate}[label=\textbf{\arabic*.},start=1]
\item \textbf{Распределение сотрудников.} В компании работают $20$ сотрудников, которых необходимо направить на два объекта. Если на первом объекте заняты $t$ человек, общая дневная оплата их труда составляет $2t^2$ условных единиц. Если на втором объекте заняты $u$ человек, общая оплата их труда составляет $u^2$ условных единиц. Все сотрудники должны быть распределены между объектами. Сколько сотрудников следует направить на каждый объект, чтобы суммарная оплата труда была наименьшей? Чему равна эта оплата?

\item \textbf{Две строительные площадки.} Бригада из $30$ рабочих должна быть полностью распределена между двумя строительными площадками. При работе $t$ человек на первой площадке суточные расходы на оплату труда равны $4t^2$ условных единиц. При работе $u$ человек на второй площадке расходы равны $u^2$ условных единиц. Найдите распределение рабочих, при котором суммарные расходы будут минимальны, и определите их величину.

\item \textbf{Рабочие в двух цехах.} На фабрике имеется $18$ рабочих. Всех рабочих распределяют между двумя цехами. Если в первом цехе работает $t$ человек, их общая дневная зарплата составляет $9t^2$ рублей; если во втором работает $u$ человек, их общая зарплата составляет $u^2$ рублей. Какое количество рабочих необходимо направить в каждый цех, чтобы сумма зарплат была минимальной? Укажите эту сумму.
\end{enumerate}

\subsection*{Средний уровень}
\begin{enumerate}[label=\textbf{\arabic*.},start=4]
\item \textbf{Заказ на электронные приборы.} Два завода получили общий заказ на изготовление $15$ электронных приборов за неделю. На первом заводе изготовление $t_1$ приборов требует $4t_1^3$ человеко-часов, а на втором изготовление $t_2$ приборов требует $t_2^3$ человеко-часов. Оба завода оплачивают каждый человеко-час по ставке $1000$ рублей. Как распределить заказ, чтобы суммарные затраты на оплату труда были минимальны? Укажите минимальную сумму в рублях.

\item \textbf{Производство деталей.} Два предприятия вместе должны изготовить $14$ деталей. На первом предприятии при выпуске $t_1$ деталей трудозатраты составляют $2t_1^3$ человеко-часов, а на втором при выпуске $t_2$ деталей~--- $t_2^3$ человеко-часов. Стоимость каждого человеко-часа составляет $1000$ рублей на обоих предприятиях. Сколько деталей следует изготовить на каждом предприятии для минимизации расходов на оплату труда? Какова минимальная сумма расходов?

\item \textbf{Распределение еженедельного заказа.} Две производственные линии должны совместно изготовить ровно $21$ изделие. На первой линии выпуск $t_1$ изделий требует $9t_1^3$ человеко-часов, а на второй выпуск $t_2$ изделий~--- $t_2^3$ человеко-часов. Один человеко-час оплачивается по ставке $1000$ рублей. Найдите целочисленное распределение выпуска, при котором суммарные затраты на оплату труда будут минимальны. Определите эти затраты в рублях.
\end{enumerate}

\subsection*{Уровень выше среднего}
\begin{enumerate}[label=\textbf{\arabic*.},start=7]
\item \textbf{Минимальная цена товара.} Предприятие ежегодно выпускает $x$ тысяч единиц товара и продаёт каждую единицу по цене $p$ тысяч рублей. Ежегодная выручка равна $px$ млн рублей, а ежегодные расходы (в млн рублей) задаются формулой $q(x)=x^2+6x+9$. План выпуска и цена не меняются в течение четырёх лет. Какова наименьшая возможная цена $p$, при которой суммарная прибыль за четыре года составит не менее $160$ млн рублей? Сколько тысяч единиц товара следует выпускать ежегодно? Считайте $x>0$.

\item \textbf{Выбор цены на продукцию.} Компания планирует ежегодно выпускать $x$ тысяч одинаковых изделий и продавать каждое по цене $p$ тысяч рублей. Годовая выручка равна $px$ млн рублей; годовые расходы (в млн рублей) равны $q(x)=0{,}5x^2+2x+10$. Объём выпуска и цена остаются постоянными три года. Определите минимальную цену $p$, при которой суммарная прибыль за три года будет не меньше $66$ млн рублей. Найдите соответствующий ежегодный выпуск. Считайте $x>0$.
\end{enumerate}

\subsection*{Сложный уровень}
\begin{enumerate}[label=\textbf{\arabic*.},start=9]
\item \textbf{Максимальный выпуск при ограниченном бюджете.} Два завода производят одинаковые изделия. Если первый завод выпускает $t_1$ изделий за неделю, зарплата его работников составляет $160t_1^2$ рублей. Если второй завод выпускает $t_2$ изделий за неделю, зарплата его работников составляет $360t_2^2$ рублей. На оплату труда обоих заводов выделено не более $74\,880$ рублей в неделю. Какое наибольшее общее число изделий можно произвести за неделю? Укажите объём выпуска каждого завода. Число изделий на каждом заводе должно быть целым неотрицательным.

\item \textbf{Оптимизация производственного плана.} Компания владеет двумя предприятиями, выпускающими один вид продукции. При производстве $t_1$ изделий на первом предприятии недельные расходы на оплату труда составляют $180t_1^2$ рублей, а при производстве $t_2$ изделий на втором~--- $320t_2^2$ рублей. Общий недельный бюджет на оплату труда ограничен суммой $288\,000$ рублей. Сколько изделий максимально можно изготовить суммарно за неделю и как распределить их выпуск между предприятиями? Оба объёма производства должны быть целыми неотрицательными числами.
\end{enumerate}

\endgroup
\clearpage
\section*{Ответы к тренировочному блоку}
\renewcommand{\arraystretch}{1.48}
\begin{longtable}{@{}p{9mm}p{45mm}p{90mm}@{}}
\toprule
\textbf{№} & \textbf{Оптимальное распределение / выпуск} & \textbf{Ответ}\\
\midrule
\endhead
1 & $t=7$, $u=13$ & Минимальная оплата: $267$ у.~е.\\
2 & $t=6$, $u=24$ & Минимальные расходы: $720$ у.~е.\\
3 & $t=2$, $u=16$ & Минимальная зарплата: $292$ руб.\\
4 & $t_1=5$, $t_2=10$ & Минимальная сумма: $1\,500\,000$ руб.\\
5 & $t_1=6$, $t_2=8$ & Минимальная сумма: $944\,000$ руб.\\
6 & $t_1=5$, $t_2=16$ & Минимальная сумма: $5\,221\,000$ руб.\\
7 & $x=7$ тыс. изделий в год & $p_{\min}=20$ тыс. руб. за изделие.\\
8 & $x=8$ тыс. изделий в год & $p_{\min}=10$ тыс. руб. за изделие.\\
9 & $t_1=18$, $t_2=8$ & Максимальный выпуск: $26$ изделий.\\
10 & $t_1=32$, $t_2=18$ & Максимальный выпуск: $50$ изделий.\\
\bottomrule
\end{longtable}

\section*{Итоговая памятка}
\begin{tabularx}{\linewidth}{@{}p{0.36\linewidth}X@{}}
\toprule
\textbf{Сигнал в условии} & \textbf{Что сделать}\\
\midrule
«Всего $n$ работников/изделий» & Записать $x+y=n$, выразить $y=n-x$.\\
«Если выпускают $x$...» & Использовать указанную формулу как функцию от числа выпускаемых изделий.\\
«Не более $B$ рублей» & Записать бюджетное ограничение $C_1+C_2\le B$.\\
«Наименьшая цена $p$» & Выразить нижнюю допустимую границу $p$ через выпуск и минимизировать её.\\
«Наибольшее количество» & Максимизировать число изделий; зарплата может играть роль ограничения.\\
«Рабочих / изделий целое число» & Проверить допустимые целочисленные варианты и границы.\\
\bottomrule
\end{tabularx}

\vfill
\begin{center}
\footnotesize\color{textgray} Лёвин Артём Александрович \quad\textbullet\quad эксклюзивно для Матвея Горбачева
\end{center}
\end{document}
